\documentclass[10pt]{article}
\usepackage[a4paper, margin=1cm]{geometry}\usepackage[brazilian]{babel}
\usepackage{enumitem}
\usepackage{xcolor}
\usepackage{tcolorbox}
\usepackage{qrcode}
\usepackage{fontawesome5}
\pagestyle{empty}

\begin{document}

\section*{\faBitcoin~Gerando uma Frase Semente BIP39 com Dados de 20 Faces}

\footnotesize

Uma frase semente mnemônica BIP39 fornece a entropia mestre para sua carteira de Bitcoin. Embora softwares possam gerar essa aleatoriedade, o uso de dados físicos oferece uma fonte de entropia transparente e verificável que não depende de nenhum gerador de números aleatórios eletrônico. Este tutorial explica como gerar uma frase semente segura de 12 palavras utilizando um dado de 20 faces e a tabela de consulta (verso).

\subsection*{\faTable~Entendendo a Estrutura da Tabela}

A tabela de referência está organizada como um sistema de consulta 
tridimensional. Cada uma das 2048 palavras do padrão BIP39 é identificada de forma única por três coordenadas, que chamaremos de D1, D2 e D3. A primeira coordenada (\textbf{D1}) seleciona uma das \textbf{8 seções} de página, a segunda (\textbf{D2}) seleciona uma das \textbf{16 linhas} dentro dessa seção, e a terceira (\textbf{D3}) seleciona uma das \textbf{16 colunas}. Juntos, esses três valores identificam exatamente uma palavra na tabela.

\subsection*{\faDiceD20~O Procedimento de Rolagem}

Para cada uma das primeiras 11 palavras da sua frase semente, você deve obter três resultados válidos do dado. Role o dado e registre o resultado apenas se ele mostrar um valor entre 1 e 16 (inclusive). Se o dado mostrar 17, 18, 19 ou 20, \textbf{descarte essa rolagem} e role novamente até obter um resultado válido. Repita esse processo até ter três números válidos (D1, D2, D3) para a palavra atual.

\begin{tcolorbox}[colback=yellow!15, colframe=yellow!65!black, title=\faLightbulb~Exemplo de Geração de Palavra]
Suponha que você esteja gerando sua quinta palavra da semente. Você rola o dado e obtém: 19 (inválido, role novamente), 3 (válido, D1 = \texttt{3,4}), 1 (válido, D2 = \texttt{1}), 20 (inválido, role novamente), 11 (válido, D3 = \texttt{11}). Suas coordenadas são ((3,4), 1, 11). Navegue até a seção 3,4 na tabela, encontre a linha 1 e então a coluna 11 para obter sua palavra: \texttt{candy}.
\end{tcolorbox}

\subsection*{\faPen~12ª Palavra}

Para gerar a última palavra da sua semente BIP39, repita o procedimento para as 2 primeiras coordenadas igual ao realizado para as 11 primeiras palavras, selecionando, portanto, a seção e a linha da tabela. Neste ponto, existe exatamente uma única palavra na linha selecionada (com 16 alternativas) que formará uma semente BIP39 válida. A determinação da coluna para a última palavra \textbf{não} será feita por rolamento de dados, e sim por tentativa e erro no wizard de input da semente da sua carteira.

\begin{tcolorbox}[colback=yellow!15, colframe=yellow!65!black, title=\faLightbulb~`Exemplo de Geração da Última Palavra]
Suponha que suas 11 primeiras palavras sejam, em ordem: \textit{``never, use, this, example, because, private, key, secret, need, phrase,} e \textit{true''}.\\

Agora você gerará sua 12ª palavra da semente. Você rola o dado e obtém: 18 (inválido, role novamente), 12 (válido, D1 = \texttt{11,12}), 9 (válido, D2 = \texttt{9}). Com isso, a sua 12ª palavra já está determinada: teste uma por uma as palavras da linha \texttt{9} da seção \texttt{11,12} e você verá que palavra da coluna \texttt{14}, \texttt{random}, (e apenas ela), formará, juntamente com as 11 primeiras, nesta ordem, uma semente BIP39 válida: 

\vspace{1em}
\texttt{1.never 2.use 3.this 4.example 5.because 6.private 7.key 8.secret 9.need 10.phrase 11.true 12.random}. \faCheckCircle
\vspace{1em}
\end{tcolorbox}

\subsection*{\faListOl~Resumo}

\begin{itemize}
 \item Role um dado de 20 faces. \textbf{Descarte} resultados \textbf{17, 18, 19} ou \textbf{20}
 \item Cada rolamento válido (de 1 a 16 inclusive) especifica uma coordenada de uma palavra
 \item Cada uma das \textbf{8 seções (D1)} corresponde a \textbf{2} resultados válidos possíveis
 \item Cada uma das \textbf{16 linhas (D2)}, bem como cada uma das \textbf{16 colunas (D3)} corresponde a \textbf{1} único resultado válido possível
 \item A \textbf{coluna (D3) da 12ª palavra}, excepcionalmente, é determinada por tentativa e erro de uma em \textbf{16} possibilidades
\end{itemize}

\vspace{1em}
\begin{tcolorbox}[colback=yellow!10, colframe=orange!90, title=\faExclamationTriangle\ Avisos de Segurança]
\begin{itemize}
 \item Nunca use a frase semente deste exemplo acima, porque para as suas chaves privadas serem secretas, a sua frase semente precisa ser verdadeiramente randômica
 \item Realize este procedimento em um local privado. Não fotografe nem registre digitalmente suas rolagens ou resultados intermediários. Escreva sua frase semente final em papel e guarde-a em local seguro
 \item Para máxima segurança, apenas digite sua frase semente em dispositivos \textit{air gapped} confiados e de sua propriedade
\end{itemize}

\end{tcolorbox}

\vspace{1em}

\begin{tcolorbox}[colback=blue!5, colframe=blue!50, title=\faInfoCircle~Saiba Mais]
\begin{minipage}[c]{0.8\textwidth}
Gostou deste tutorial? Acesse: \texttt{www.loudproudandfree.com} pelo seu navegador ou escaneie o QR code ao lado. Lá você encontrará:

\begin{description}
 \item [Mais tutoriais:] Gostaria de \textbf{frases com tema} em vez de palavras soltas? Temos!
 \item [Cursos e consultorias:] Tire suas dúvidas com um especialista
 \item [Twitter, Nostr e Telegram:] Fala parte do movimento e siga-nos para se manter informado
 \item [Segurança Física:] Medo de roubo ou sequestro? Confira nosso app de segurança física anti-frágil
\end{description}

\end{minipage}%
\hfill
\begin{minipage}[c]{0.15\textwidth}
\centering
\qrcode[height=2.5cm]{https://www.loudproudandfree.com}

% \small\textit{Escaneie para acessar}
\end{minipage}
\end{tcolorbox}


\end{document}
